\infowidth{8cm}

\cheading{天津大学  机器学习与深度学习课程实验报告}
\ctitle{基于多源海洋监测数据的盐度变化规律分析与预测}
\caffil{人工智能学院}
\csubject{人工智能}
\cgrade{24级}
\cauthor{卢思宇}
\cnumber{3024244246}
\csupervisor{王煜}

\cabstractcn{
本实验面向真实海洋环境监测场景，围绕“盐度变化规律分析与预测”这一核心问题，使用海流、波浪、水质、气象与设备状态五类异构监测数据，构建了多源数据融合、分颗粒度建模与结果解释的一体化分析流程。针对真实数据中普遍存在的异步采样、缺失值与异常值问题，本文以水质表为锚点，按\texttt{ZDMC + DataTime}建立多表左连接，结合物理范围约束、分类型缺失处理与时间周期特征编码，形成可用于监督与无监督任务的统一样本。建模策略遵循“先小颗粒度、再大颗粒度、最后结构解释”的路径：先在小时级建立线性回归基线，再在天级聚合数据上进行对比建模，并使用时间序列交叉验证选择Ridge正则强度，最后通过KMeans聚类识别环境状态类型并分析其物理特征差异。\\
\hspace*{\parindent}实验结果表明，在严格的“历史训练—未来测试”设置下，Ridge回归在小时级与天级均优于普通线性回归，小时级最佳模型取得\texttt{RMSE=8.793}、\texttt{MAE=3.669}，天级最佳模型取得\texttt{RMSE=12.738}、\texttt{MAE=6.013}，说明正则化在线性体系内对多源高噪声数据具有显著稳健化作用。聚类结果显示最优簇数为2（\texttt{silhouette=0.277}），尽管簇间分离度中等，但仍揭示了样本在盐度、浊度、风速和流速上的结构差异，为回归误差来源解释提供了辅助证据。综合来看，本实验在课程已学方法框架内实现了从数据治理到模型评估、从数值预测到规律解释的完整闭环，结果具有可复现性与分析价值。
}

\ckeywordcn{海洋环境监测，多源数据融合，盐度预测，Ridge回归，时间颗粒度，聚类分析}

\cabstracten{}
\ckeyworden{}
